% Generated draft: requires mathematical and source review.
\documentclass[11pt]{article}
\usepackage[margin=25mm]{geometry}
\usepackage{fontspec}
\setmainfont{Latin Modern Roman}
\usepackage{amsmath,amssymb,mathtools,unicode-math}
\setmathfont{Latin Modern Math}
\usepackage{enumitem,needspace,xurl,hyperref}
\setlength{\parindent}{0pt}
\setlength{\parskip}{5pt}
\setlength{\emergencystretch}{3em}
\newcommand{\N}{\mathbb{N}}
\newcommand{\Z}{\mathbb{Z}}
\newcommand{\Q}{\mathbb{Q}}
\newcommand{\R}{\mathbb{R}}
\newcommand{\C}{\mathbb{C}}
\newcommand{\F}{\mathbb{F}}
\newcommand{\A}{\mathbb{A}}
\newcommand{\PP}{\mathbb{P}}
\newcommand{\E}{\mathbb{E}}
\newcommand{\eps}{\varepsilon}
\newcommand{\dd}{\,\mathrm{d}}
\begin{document}
\section*{UnsolvedMath ID: 30003547}
\textbf{Source problem code:} catalogue entry
\subsection{Definitions and mathematical statement}
Let \(L\) be a Lie manifold equipped with a fixed compatible positive-scalar-curvature metric \(g_0\), in the sense intended by the supplied source. Positive scalar curvature means that \(\operatorname{Scal}_{g_0}(x)>0\) at every point where the associated Riemannian metric is defined. The record does not specify the definition of a Lie manifold, metric compatibility, or the requirements at boundary faces or at infinity.

Write \(\mathcal C(L;g_0)\) provisionally for the intended collection of concordance classes of metrics relative to \(g_0\). This notation is a placeholder: the admissible metrics and the concordance relation must be specified before it defines a mathematical object.

\textbf{Source question.} In the category of Lie manifolds, can one define an appropriate composition law on concordance classes of metrics relative to a fixed compatible positive-scalar-curvature metric?

If composition is intended to be a total binary operation, the requested construction would have the form
\[
\mu_{L,g_0}\colon \mathcal C(L;g_0)\times\mathcal C(L;g_0)\longrightarrow\mathcal C(L;g_0),
\]
with its value independent of representatives. The source does not determine whether composition is total or partially defined, what algebraic axioms it must satisfy, or what compatibility with morphisms of Lie manifolds is required. In particular, it does not explicitly require a group structure or identify \([g_0]\) as an identity. The meaning of ``relative to'' and the precise quantification over Lie manifolds are also unresolved. Consequently, this is a formulation of the supplied structural question with explicit gaps, rather than a fully specified existence conjecture.
\subsection{Short English statement}
The question asks how to compose concordance classes of metrics on a Lie manifold after fixing a compatible positive-scalar-curvature metric. The supplied record leaves the concordance relation, relative conditions, and required properties of composition unspecified, so a precise existence statement requires further source clarification.
\subsection{Sources}
\begin{itemize}
\item Local UnsolvedMath registry, supplied complete source record: permanent ID \(30003547\), title \emph{Composition of Positive-Scalar-Curvature Concordance Classes}, source problem code \texttt{catalogue entry}. The record assigns proposed year 2017 and names no individual proposer.
\item Supplied report citation: \emph{Analysis, Geometry and Topology of Positive Scalar Curvature Metrics} (2018), Oberwolfach Reports, pp.~2223--2298, DOI \texttt{10.4171/owr/2017/36}. Original OWR report: \url{https://doi.org/10.4171/owr/2017/36}. The supplied background describes the report as documenting workshop discussions and open problems.
\item Supplied literature reference, labeled ``relative higher index and PSC concordance'': \url{https://arxiv.org/abs/1712.03722}. No authors, bibliographic title, or specific result are supplied for this reference.
\item Literature assessment dated exactly \texttt{2026-08-21}: the supplied curation labels the problem open and reports that group structures on concordance classes exist in important closed-manifold settings after fixing a positive-scalar-curvature metric, while the requested composition for Lie manifolds remains an open structural problem. This is supplied dated curation derived from the OWR material, not an independently verified literature assessment. The supplied background expressly requests independent verification of time-sensitive status and cited claims.
\end{itemize}
\end{document}
